\section{Design Lessons}
\label{sec:disc-lessons}

Correct interpretation of the half-bridge current input PA2 was essential
to the control law. Secondary voltage multiplied by this signal does not
represent output power; direct secondary-voltage feedback was therefore
used, with PA2 retained for current sensing.

Acquisition latency, enable-edge alignment, and minimum packet length
must be designed with the tank, output capacitance, and sensing bandwidth.
Separate no-load limits, holdoff, and recovery logic were needed to reduce
charging after disconnection.

The V2 input hotspot persisted after diode replacement, so the current and
thermal paths around the bridge need further investigation. One possible
restriction is a single-via connection between layers: narrow pad connections
and interlayer paths can restrict current and heat flow despite a large
copper pour.

Secondary-winding heating motivated transformer optimization, but a
replacement set was not installed. The V3 layout reduced volume by removing
the USB-PD section, so a complete-system power-density comparison must also
include the equivalent external output functions.

Recording the board version, settings, waveforms, power readings, and
temperatures together allows each result to be interpreted against its test
conditions. Coupling, averaging, and firmware state are also needed when
calculating results from captures.
